\documentclass[11pt]{article}

\input{0-parameters.tex}

\begin{document}

\section{Combining Assumptions}


We take as accounting truth that
\[
    \hat{A} = \sum_{i=1}^N \lambda_i \hat{A}_i,
\]
and we define
\[
\hat{\tilde{A}}_i \equiv \left(\zeta\hat{X}_i+(\zeta - 1) \hat{p}_i\right).
\]

We also know that, by definition of $\zeta_i$,
\[
\hat{A}_i \equiv \left(\zeta_i \hat{X}_i+(\zeta_i - 1) \hat{p}_i\right).
\]

In other words, we can write
\begin{equation}
\hat{A}   = \sum_{i=1}^N \lambda_i \hat{\tilde{A}}_i + \sum_{i=1}^N \lambda_i \mathcal{E}_i,
\label{eq:ahat_plus_error}
\end{equation}
where
\begin{equation}
\mathcal{E}_i \equiv  (\zeta_i - \zeta) \left( \hat{X}_i + \hat{p}_i\right).
\label{eq:eps}
\end{equation}


Now, we show (see below section \ref{section_production_network_approach}) that 
\begin{equation}
\label{eq_agg_A_appendix}
\sum_{i=1}^N \lambda_i \hat{\tilde{A}}_i = \sum_{f=1}^F \frac{w_fL_{f}}{pY}\hat{w}_f,
\end{equation}
which happens to be the dual definition of aggregate TFP\footnote{
Following \cite{hsieh_productivity_1999} and \cite{antras_factor_2003}, we start from the accounting identity that total value added, GDP, $Y$ is equal to the sum of factor payments
\begin{equation*}
    Y=\sum_{f=1}^F w_f L_f.
\end{equation*}
Differentiating, we get
\begin{equation*}
    \hat{Y}=\sum_{f=1}^F \frac{w_fL_{f}}{pY}(\hat{L}_f+\hat{w}_f),
\end{equation*}
which we rearrange as
\begin{equation*}
    \underbrace{\hat{Y}-\sum_{f=1}^F \frac{w_fL_{f}}{pY}\hat{L}_f}_{\text{Agg. TFP primal}}= \underbrace{\sum_{f=1}^F \frac{w_fL_{f}}{pY}\hat{w}_f}_{\text{Agg. TFP dual}},
\end{equation*}
where the left-hand side is the classic definition of the aggregate Solow residual in its primal form, thus equal to the weighted sum of the growth rate of real factor prices, so-called dual.

},
\begin{equation}
\label{eq_agg_A_dual}
\hat{A} = \sum_{f=1}^F \frac{w_fL_{f}}{pY}\hat{w}_f.
\end{equation}

Inserting Eq. \eqref{eq_agg_A_appendix} and Eq. \eqref{eq_agg_A_dual} into Eq. \eqref{eq:ahat_plus_error}, we get
\begin{equation}
\label{eq_weighted_sum_error}
    \sum_{i=1}^N \lambda_i \mathcal{E}_i = 0.
\end{equation}

In other words, one could use Equation~\ref{eq_weighted_sum_error} to calibrate a universal value of $\zeta$ by minimising the aggregate error—specifically, the Domar-weighted sum of industry-level errors. This would directly reduce the overall productivity mismeasurement. However, our approach (see Section~\ER{XXX}) instead minimises the \emph{unweighted} sum of sector-level errors and then evaluates the aggregation error in a second step. This implicitly assigns more weight to smaller sectors (in terms of Domar weights). We prefer this strategy for two reasons:
(i) it ensures that both steps—sector-level estimation and aggregation—yield unbiased and well-controlled errors;
(ii) since the net-zero transition is unlikely to significantly affect the largest sectors (such as services), we avoid over-weighting them in the calibration process.

\subsubsection{Production network approach -- Proof of Eq. \eqref{eq_agg_A_appendix}}
\label{section_production_network_approach}

We have N sectors in the economy. Hence, industry-level productivity prediction is $\hat{\tilde{A}}_i = \zeta\hat{X}_i+(\zeta - 1) \hat{p}_i$.
Let us concatenate the two equations, assuming that $\hat{A}_i\approx \hat{\tilde{A}}_i$ (which we prove numerically). Let us remember that the dual definition of productivity gives $\hat{A}_i \equiv  - \left( \hat{p}_i - \sum_{j=1}^N  a_{ji} \hat{p}_j - \sum_{f=1}^F \tilde{l}_{fi} \hat{w}_f \right)$.
Hence,
\begin{equation*}
     \hat{p}_i=-\hat{A}_i + \sum_{j=1}^N  a_{ji} \hat{p}_j+\sum_{f=1}^F \tilde{l}_{fi} \hat{w}_f.
\end{equation*}
And the primal, $\hat{A}_i \equiv \hat{X}_i - \sum_{j=1}^N  a_{ji} \hat{X}_{ij} - \sum_{f=1}^F \tilde{l}_{fi}\hat{L}_{fi}$, then
\begin{equation*}
     \hat{X}_i=\hat{A}_i + \sum_{j=1}^N  a_{ji} \hat{X}_{ij} + \sum_{f=1}^F \tilde{l}_{fi}\hat{L}_{fi}.
\end{equation*}

\begin{align*}
    \hat{A}&=\sum_{i=1}^N \lambda_i \hat{A}_i\\
    & \approx \sum_{i=1}^N \lambda_i \hat{\tilde{A}}_i\\
    & = \sum_{i=1}^N \lambda_i \left(\zeta\hat{X}_i+(\zeta - 1) \hat{p}_i\right) \\
     & = \zeta\sum_{i=1}^N \lambda_i\hat{X}_i+(\zeta - 1) \sum_{i=1}^N \lambda_i\hat{p}_i \\
      & = \zeta\sum_{i=1}^N \lambda_i \left(\hat{A}_i + \sum_{j=1}^N  a_{ji} \hat{X}_{ij} + \sum_{f=1}^F \tilde{l}_{fi}\hat{L}_{fi} \right)+(\zeta - 1) \sum_{i=1}^N \lambda_i\left(-\hat{A}_i + \sum_{j=1}^N  a_{ji} \hat{p}_j+\sum_{f=1}^F \tilde{l}_{fi} \hat{w}_f\right) \\
      & = \sum_{i=1}^N \lambda_i\hat{A}_i
      +\zeta \underbrace{\left(
      \sum_{i=1}^N \sum_{j=1}^N \lambda_i a_{ji} \hat{X}_{ij} + \sum_{i=1}^N\sum_{f=1}^F  \lambda_i\tilde{l}_{fi}\hat{L}_{fi}
      \right)}_{A}+
      (\zeta - 1) \underbrace{\left(
      \sum_{i=1}^N \sum_{j=1}^N   \lambda_ia_{ji} \hat{p}_j+\sum_{i=1}^N \sum_{f=1}^F  \lambda_i\tilde{l}_{fi} \hat{w}_f
      \right) }_{B}
\end{align*}

We simplify this equation, using the following
\begin{equation}
    \begin{cases}
\hat{A} \equiv \sum_{i=1}^N \frac{p_iC_i}{pY} \hat{C}_i - \sum_{f=1}^F \frac{w_f L_f}{pY} \hat{L}_f,\\
 L_f \hat{L}_f = \sum_i  L_{fi} \hat{L}_{fi},\\
    X_j=\sum_i X_{ij} + C_j,\\
    \sum_i X_{ij}\hat{X}_{ij} = X_j (\hat{X}_j - (C_j/X_j) \hat{C}_j).
    \end{cases}
\end{equation}



We can simplify term A,
\begin{align*}
\left(
      \sum_{i=1}^N \sum_{j=1}^N \lambda_i a_{ji} \hat{X}_{ij} + \sum_{i=1}^N\sum_{f=1}^F  \lambda_i\tilde{l}_{fi}\hat{L}_{fi}
      \right)&=\left(
      \sum_{i=1}^N \sum_{j=1}^N \frac{p_iX_i}{pY} \frac{p_jX_{ij}}{p_iX_i} \hat{X}_{ij} + \sum_{i=1}^N\sum_{f=1}^F  \frac{p_iX_i}{pY} \frac{w_fL_{fi}}{p_iX_i}\hat{L}_{fi}
      \right)\\
      &=
     \sum_{i=1}^N \sum_{j=1}^N \frac{p_jX_{ij}}{pY} \hat{X}_{ij} + \sum_{i=1}^N\sum_{f=1}^F  \frac{w_fL_{fi}}{pY}\hat{L}_{fi}\\
     &=
      \sum_{j=1}^N\frac{p_j}{pY}\sum_{i=1}^N   X_{ij}\hat{X}_{ij} +\sum_{f=1}^F  \frac{w_f}{pY}  \sum_{i=1}^N L_{fi}\hat{L}_{fi}\\
      &=
      \sum_{j=1}^N\frac{p_j}{pY}\left(X_j (\hat{X}_j - (C_j/X_j) \hat{C}_j)\right) +\sum_{f=1}^F  \frac{w_f}{pY}   L_f \hat{L}_f\\
      &=
      \sum_{j=1}^N\lambda_j\hat{X}_j - \sum_{j=1}^N\frac{p_jC_j}{pY}  \hat{C}_j +\sum_{f=1}^F  \frac{w_f L_f }{pY}  \hat{L}_f.\\
\end{align*}


We can simplify term B,
\begin{align*} 
    \left(
      \sum_{i=1}^N \sum_{j=1}^N  \frac{p_iX_i}{pY} \frac{p_jX_{ij}}{p_iX_i} \hat{p}_j+\sum_{i=1}^N \sum_{f=1}^F  \frac{p_iX_i}{pY} \frac{w_fL_{fi}}{p_iX_i} \hat{w}_f
      \right)
      &=\sum_{i=1}^N \sum_{j=1}^N   \frac{p_jX_{ij}}{pY} \hat{p}_j+\sum_{i=1}^N \sum_{f=1}^F \frac{w_fL_{fi}}{pY} \hat{w}_f\\
      &=
       \sum_{j=1}^N\hat{p}_j \frac{p_j}{pY}\sum_{i=1}^N   X_{ij} + \sum_{f=1}^F \hat{w}_f\frac{w_f}{pY}\sum_{i=1}^N L_{fi} \\
       &=
       \sum_{j=1}^N\hat{p}_j \frac{p_j}{pY}X_{j} + \sum_{f=1}^F \hat{w}_f\frac{w_f}{pY}L_{f} \\
        &=
       \sum_{j=1}^N\lambda_j \hat{p}_j + \sum_{f=1}^F \frac{w_fL_{f}}{pY}\hat{w}_f \\
\end{align*}


Combining all equations back together, 

\begin{align*}
    \hat{A}&  = \sum_{i=1}^N \lambda_i\hat{A}_i
      +\zeta \left(
      \sum_{j=1}^N\lambda_j\hat{X}_j - \sum_{j=1}^N\frac{p_jC_j}{pY}  \hat{C}_j +\sum_{f=1}^F  \frac{w_f L_f }{pY}  \hat{L}_f
      \right)+
      (\zeta - 1) \left(
       \sum_{j=1}^N\hat{p}_j \lambda_j + \sum_{f=1}^F \frac{w_fL_{f}}{pY}\hat{w}_f
      \right) \\
       & = \sum_{i=1}^N \lambda_i\hat{A}_i
      + \sum_{j=1}^N\lambda_j(\zeta \hat{X}_j+ (\zeta - 1)\hat{p}_j)
      +\zeta \left(
     - \sum_{j=1}^N\frac{p_jC_j}{pY}  \hat{C}_j +\sum_{f=1}^F  \frac{w_f L_f }{pY}  \hat{L}_f
      \right)+
      (\zeta - 1) \left(
       \sum_{f=1}^F \frac{w_fL_{f}}{pY}\hat{w}_f
      \right) \\
       & = \sum_{i=1}^N \lambda_i\hat{A}_i
      + \sum_{j=1}^N\lambda_j\hat{\tilde{A}}_j
      +\zeta \left(
    -\hat{A}
      \right)+
      (\zeta - 1) \left(
       \sum_{f=1}^F \frac{w_fL_{f}}{pY}\hat{w}_f
      \right) \\
      \Rightarrow & \hat{A}(1-2+\zeta)= 
      (\zeta - 1) \left(
       \sum_{f=1}^F \frac{w_fL_{f}}{pY}\hat{w}_f
      \right) \\
      \Rightarrow & \hat{A}(\zeta-1)= 
      (\zeta - 1) \left(
       \sum_{f=1}^F \frac{w_fL_{f}}{pY}\hat{w}_f
      \right).
\end{align*}

Hence,
either
\begin{equation}
    \begin{cases}
        \zeta=1\\
        \textbf{ or } \\
        \hat{A}=  \sum_{f=1}^F \frac{w_fL_{f}}{pY}\hat{w}_f.
    \end{cases}
\end{equation}



% Bibliography ---------------------------
\small
\renewcommand\bibname{References}
\bibliographystyle{apalike}
\bibliography{z-bib}
\normalsize

\end{document}